\section{The centred-CRT control system}\label{sec:control}

Fix a squarefree integer $b\ge3$ and a finite obstruction set $T$.  The bottom
scale $k_0$ will be chosen only after all other parameters.  Put $K=3k_0$ and
use the full dyadic blocks
\[
 P_k=\Pk_k,
 \qquad
 |P_k|\ge\frac{2^k}{2\log(2^k)}
 \quad(k_0\le k\le3k_0).
\]
The estimates remain valid for subsets with the same lower density, but the
full blocks ensure that the later mass family and prime reservoir belong to one
period support.  The half-open intervals are disjoint, and
\[
 \Bset=\bigcup_{k=k_0}^{K}P_k.
\]
We take all block primes outside the prime support of $b$ and large enough that
the products used below avoid $T$.

The control graph consists of every unordered pair inside one block and every
pair between consecutive blocks:
\begin{equation}\label{eq:control-graph}
 \Cgraph=
 \bigcup_{k=k_0}^{K}\{(p,q):p,q\in P_k,\ p<q\}
 \;\cup\;
 \bigcup_{k=k_0}^{K-1}(P_k\times P_{k+1}).
\end{equation}
The two indexed families are disjoint.  Unique factorization makes the product
map injective on them, so
\[
 \Ectrl=\{pq:(p,q)\in\Cgraph\}
\]
is a set of distinct squarefree semiprimes.

A global assignment is a residue vector
\[
 a=(a_p)_{p\in\Bset},
 \qquad a_p\in\Z/p\Z.
\]
For distinct block primes $p,q$, let $H_{pq}(a)$ be the unique centred CRT
representative satisfying
\begin{equation}\label{eq:centered-crt}
 H_{pq}(a)\equiv a_p\pmod p,
 \qquad
 H_{pq}(a)\equiv a_q\pmod q,
 \qquad
 -\frac{pq}{2}<H_{pq}(a)\le\frac{pq}{2}.
\end{equation}
Define
\begin{equation}\label{eq:Qctrl}
 \Qctrl(a)=\sum_{(p,q)\in\Cgraph}
 \left(\frac{H_{pq}(a)}{pq}\right)^2,
 \qquad
 \sctrl^2=\sum_{(p,q)\in\Cgraph}\frac1{(pq)^2}.
\end{equation}
For a Fourier frequency $h$, let $a(h)_p=h\pmod p$.  The centred convention gives
the exact phase identity
\begin{equation}\label{eq:control-energy-identity}
 \normdist{h/(pq)}=\frac{|H_{pq}(a(h))|}{pq},
 \qquad
 \Qctrl(a(h))\le\QE(h):=\sum_{e\in E}\normdist{h/e}^2.
\end{equation}
This centring is also what later turns a globally constant residue assignment
into an exact one-dimensional quadratic energy.

The control mechanism needs two scale estimates.  The first requires an upper
bound for the number of primes in a block; large denominators alone do not
control the number of edges.

\begin{lemma}[Elementary block reciprocal bound]\label{lem:block-upper}
There is an absolute constant $C_B>0$ such that, for every subset $P_k$ of the
primes in $[2^k,2^{k+1})$ and every $k\ge1$,
\begin{equation}\label{eq:block-recip-upper}
 \sum_{p\in P_k}\frac1p\le\frac{C_B}{k}.
\end{equation}
\end{lemma}
\begin{proof}
Let $M_k$ be the number of primes in the interval.  An elementary induction
using the central binomial coefficient gives
\[
 \prod_{p\le2^r}p\le4^{2^r}.
\]
Every prime in the $k$th block is at least $2^k$, whereas their product is at
most the primorial at $2^{k+1}$.  Hence $kM_k\le2^{k+2}$ and
$\sum_{p\in P_k}1/p\le M_k2^{-k}\le4/k$.  Thus one may take $C_B=4$.
\end{proof}

\begin{proposition}[Control load and deviation scale]\label{prop:control-estimates}
There are fixed constants $C_{\mathrm{load}},c_\sigma>0$ such that, for all
sufficiently large $k_0$,
\begin{equation}\label{eq:control-load-symbolic}
 \sum_{e\in\Ectrl}\frac1e
 \le\frac{C_{\mathrm{load}}}{k_0-1},
 \qquad
 \sctrl\ge\frac{c_\sigma}{k_0 2^{k_0}}.
\end{equation}
In particular, for fixed $b$, the control load tends to zero.
\end{proposition}
\begin{proof}
Write $S_k=\sum_{p\in P_k}1/p$.  Internal edges in block $k$ have load at most
$S_k^2$, while the adjacent bipartite load is $S_kS_{k+1}$.  Lemma
\ref{lem:block-upper} therefore gives
\[
 \sum_{e\in\Ectrl}\frac1e
 \le2C_B^2\sum_{k\ge k_0}\frac1{k^2}
 \le\frac{2C_B^2}{k_0-1}.
\]
Thus $C_{\mathrm{load}}=32$ is available from this human estimate.

For the second estimate, use only the internal edges of the bottom block.  Put
$X=2^{k_0}$ and $m=|P_{k_0}|$.  The density estimate gives
$m\gg X/k_0$, and every internal product is below $2^{2k_0+2}$.  Once $m\ge2$,
\[
 \sctrl^2
 \ge\frac{m(m-1)/2}{2^{4k_0+4}}
 \gg\frac1{k_0^2 2^{2k_0}}.
\]
Taking square roots gives the stated inequality for some fixed $c_\sigma>0$.
\end{proof}

For compatibility with the archived finite certificate, one may take
$C_{\mathrm{load}}=32$ in the preceding proof and $c_\sigma=1/100$ after a
finite threshold; the release also uses the coarser load constant $512$.  These
numbers certify \eqref{eq:control-load-symbolic} but are not independent
parameters of the construction.

For $C\ge1$, define the main assignment set
\begin{equation}\label{eq:assignment-main}
 \mathfrak M(C)=
 \left\{a:\exists m\in\Z,
 |m|\le C/\sctrl,
 \ a_p\equiv m\pmod p\text{ for all }p\in\Bset\right\}.
\end{equation}
